\section{项目预期目标}
\label{sec:project-goals}
本项目以满足课程任务要求、完成“图像描述生成—结构化膳食记录”两条主线的完整实验链路、保证全部结果可复现为总体目标，通过自主模型实现、预训练模型领域适配和固定协议下的对照实验，系统考察不同模型在食物图像理解任务上的表现。各分项目标均对应明确的交付物与检验方式。

\medskip
\noindent\textbf{目标一：食物图像描述生成（任务一，50 分）。}自主实现两种结构不同的图像描述模型：模型 A 采用 ResNet-18 全局特征经投影后驱动 GRU 循环解码，模型 B 采用 ResNet-18 网格特征经线性映射为视觉 Token 后由 Transformer 编码器—解码器生成。两种模型共享数据划分、文本规范化、训练预算与贪心解码协议，使用 BLEU-4、METEOR、ROUGE-L、CIDEr-D 等不少于三种指标在固定测试集上完成评测，并通过 B-Global 消融分析空间网格表示的贡献。检验标准为两种模型均完成从数据加载到指标输出的全流程，结果可由配置文件与 Checkpoint 复现。

\medskip
\noindent\textbf{目标二：结构化膳食记录生成（任务二，30 分）。}完成并比较两条技术路线：路线一为任务一模型 A 的参数迁移方案，采用食材序列生成、边界对齐的逐食材重量回归与整餐五维数值回归相结合的混合架构，由确定性程序组装 JSON；路线二为 Qwen3-VL-4B-Instruct 经 LoRA/QLoRA 参数高效微调直接生成 JSON。以食材名称 Micro F1、逐食材重量与整餐五维数值的 MAE/MSE、JSON 语法合法率与 Schema 有效率为指标，形成同规模原始模型、微调模型与自主混合模型的三方对照；解析失败样本按预先规定的方式计入全样本统计，保证指标口径一致。

\medskip
\noindent\textbf{目标三：任选实验（20 分）。}完成两项任选方向。其一为强化学习优化：在模型 A 交叉熵训练的基础上衔接 SCST，以 CIDEr-D 为奖励，报告优化前后四项文本指标及食材幻觉的变化，并设置等步数继续交叉熵的对照，区分额外训练与奖励优化的作用。其二为 VLM-as-Judge：以固定协议从语义正确性、细节覆盖、事实可靠性与语言流畅度四个维度评分，与双人工评分比较一致性（Spearman 相关、Bootstrap 区间与 Cohen’s $\kappa$）。

\medskip
\noindent\textbf{目标四：工程与展示。}交付统一的数据处理、训练、推理与评价代码库，配置统一记录随机种子、数据哈希、词表、数值统计与生成参数，Checkpoint 保存完整训练状态；在约 300 元、单卡 24 GB 显存的算力条件下完成全部必选实验。最终提供 FoodLens 演示原型，支持上传食物图像、生成自然语言描述、展示结构化膳食记录与模型对照结果。

